% ======================================================================
%  07_dmod.tex —— 第 7 讲 翻译器：从 A-膜到 D-模
%  依据：Kapustin--Witten 2006 \S11（印刷 pp.184--206）
% ======================================================================

\section{翻译器：从 A-膜到 D-模}
\label{sec:dmod}

\begin{goals}
\begin{itemize}
\item 说清几何 Langlands 为什么需要一台「从 A-膜到 D-模」的翻译器，以及它必须满足什么条件。
\item 掌握典范余各向同性 A-膜的定义与唯一性，并理解为什么它的开弦代数恰好是微分算子代数。
\item 记住 $f(\Psi)=\Psi$ 的结论、来源与不确定度，知道「普通 D-模」落在 $\Psi=h$ 而不是别处。
\end{itemize}
\end{goals}

几何 Langlands 的标准陈述是：曲线 $C$ 上每个局部系统 $\rho:\pi_1(C)\to\LdualC$，对应 $\MC$ 上一个既被 Hecke 算子作用、又是其本征对象的层——一个 Hecke 特征层，而 Hecke 特征层首先是一个 D-模。前几讲已经把左边翻译成了物理对象：局部系统是 magnetic eigenbrane，Hecke 算子是 't~Hooft 线算子（\cite{kw06}；\src{B p.184}）。但右边在物理侧完全缺席：我们手里只有 $\MH$ 上的膜，而 D-模住在另一个模空间上。本讲要补的就是这台翻译器，原文 \S11 把它称为「缺的一环」（\src{B p.184}）。

\subsection{缺口的形状：为什么翻译器必须「空间充满」}

要理解 \S11 的策略，先把它想成一句范畴论常识：如果手头有一张膜 $B$，那么对任意另一张膜 $B'$，$(B,B')$ 弦构成 $(B,B)$ 弦代数上的一个模；而 $(B,B)$ 弦就是这个膜作为范畴对象的自同态代数（\src{B p.196}）。于是只要把某一张膜的\textbf{自同态代数}算清楚，所有膜就自动成为这个代数的模——翻译器不是一台把膜映成 D-模的机器，而是一张\textbf{万有探针}，它的开弦代数就是答案本身。

剩下的问题是：哪张膜能当探针？普通的 A-膜支撑在 Lagrange 子流形上，维数只有靶空间的一半，它撑不满空间，也就管不到所有膜；而「所有膜」里必须包括那些支撑在别处的膜。所以探针必须满足两条：第一，它是 A-膜；第二，它撑满整个 $\MH$。前者保证它属于我们正在研究的 A-模型，后者保证它的自同态代数足够大，大到能容纳所有膜。

这里要立刻打一剂预防针。原文的承诺是「我们将会说明」每个复结构 $K$ 中的 A-膜都自动给出 $\MC$ 上的一个 D-模，措辞是 will show 加 roughly speaking（\src{B p.184}）；「范畴等价」只是转引更早的一般结论（\src{B p.184}），本情形是该结论的一个特例，原文并没有在本节证明任何范畴层面的等价。本讲把话说准：我们讲的是\textbf{单向构造}，不是函子性定理。

\subsection{典范余各向同性膜：曲率由几何唯一决定}

\term{coisotropic submanifold}{余各向同性子流形}是辛流形 $X$ 中局部由一族 Poisson 交换的函数 $f_1,\dots,f_d$ 的零点定义的子流形；这类函数最多有 $\tfrac12\dim X$ 个独立，所以 $\dim Y\ge\tfrac12\dim X$，取等号时 $Y$ 才是 Lagrange 的（\src{B p.185}）。\S11 只需要最简单的极端情形：$Y=X$、rank one，即膜是空间充满的，Chan--Paton 丛是一条 $U(1)$ 线丛 $L$，规范场局部是实一形式 $A$，曲率 $F=\dd A$。这样的膜要是 A-膜，曲率必须满足
\begin{equation}
(\omega^{-1}F)^2=-1 .
\label{eq:cc-condition}
\end{equation}
这条方程的几何含义是：$\hat N=\omega^{-1}F$ 是一个近复结构（\src{B p.185}）。由于 $F$ 与 $\omega$ 都闭，(11.1) 还保证 $\hat N$ 可积；而 $\omega$ 与 $F$ 对 $\hat N$ 都是 $(2,0)\oplus(0,2)$ 型，故 $\omega-\ii F$ 是 $\hat N$ 下的全纯辛形式。原文由此得出一个限制：这类膜只在 $\omega$ 是全纯辛形式的实部时才可能存在（\src{B p.185}）。

在 $\MH$ 上这件事恰好成立。取 $\omega=(\Im\tau)\omega_K$，则 $\hat N=\omega_K^{-1}\omega_J=I$，而 $\omega_J$ 恰当，故存在势 $\alpha$ 使 $F=(\Im\tau)\omega_J=\dd\alpha$，它是平凡线丛上的一个固定联络（\src{B p.186}）。一般解是 $F=(\Im\tau)(a\omega_J+b\omega_I)$、$a^2+b^2=1$；但当 $b\neq0$ 时该形式不恰当、周期不是 $2\pi$ 的整数倍，根本定义不了膜，所以可用的只剩 $b=0$ 与整体变号。这就是「典范」二字的含义：$F$ 不是选出来的，而是由 $\MH$ 的全纯辛几何唯一决定的；$G$ 非单连通时还可以扭一个平坦线丛，而平坦扭不改变开弦代数层（\src{B p.186}；\src{B p.193}）。

再看它的类型。$\hat N=I$ 说明它在复结构 $I$ 中是 A-膜；在 $K$ 中 $(A,B,A)$ 的判定用 $(\omega_I^{-1}\omega_J)^2=-1$ 同样成立；而在 $J$ 中 $F\propto\omega_J$ 是 $(1,1)$ 型，故它是 B-膜。三句话合起来：c.c.\ 膜是 $(A,B,A)$ 型（\src{B p.187}）。另外，它没有世界面瞬子修正：边界条件与 $K$ 中全纯性联立给出 $(\omega-\ii F)\partial_0\Phi=0$，$\omega-\ii F$ 非退化迫使 $\Phi$ 在边界上为常数，再由解析性整体为常数（\src{B p.187}）。于是计算开弦代数只需要微扰论。

开弦代数为什么是微分算子？局部看，限制到 $\MH$ 中形如 $T^*V$ 的开集（$V\subset\MC$ 小开集）后，边界态的乘法由 $\omega_J=\Re\Omega_I$ 给出的 Poisson 括号控制：$[\kappa,f]=-\ii(\Im\tau)^{-1}\kappa(f)$，而 $\hat p_\alpha=-\frac{\ii}{\Im\tau}\frac{\partial}{\partial q^\alpha}$（\src{B p.191}）。也就是说，$\varphi_z$ 方向的坐标在量子化后变成\textbf{微分}，而 $A_z$ 方向的坐标变成\textbf{乘算子}——这正是全纯微分算子的局部模型。整体上要粘：转移函数的三重上链若属于 $2\pi\ZZ$ 就得到线丛上的 $\Dmod_L$，否则得到复幂 $\Dmod_{L^\gamma}$（\src{B p.192}）。最后一步是时间反演：$t=1$ 处 $T$ 是精确对称，它把开弦代数送到 opposite 代数，而 $\Dmod_L$ 的 opposite 是 $\Dmod_{L^{-1}\otimes K_M}$；自对偶要求 $L^2=K_M$，即 $L=K_M^{1/2}$（$\MC$ 是 spin 流形，平坦扭不改层）（\src{B pp.193--195}）。这就是 c.c.\ 膜的开弦代数层 $\Dmod_{K_{\MC}^{1/2}}$，也就是几何 Langlands 里那个「所有 D-模的老家」。

\subsection{翻译如何发生：开弦代数、模与磁特征膜}

把上一小节压成一句可用的命题：任意 $B'$ 在 $K$ 中是 A-膜，则 $(B_{c.c.},B')$ 弦层构成 $\Dmod_{K_{\MC}^{1/2}}$ 的一个模——这就叫「$B'$ 给出了一个 D-模」（\src{B p.196}）。原文对此用 we claim、we argue，态度是诚实的：这是物理上自然的构造，不是数学证明。

两个例子能让人相信它。第一个例子是校准用的：取 $B'$ 为 $\varphi=0$ 处的平坦膜，它的支撑是 $\MC$ 在 $\MH$ 中的拷贝，属于 $(B,A,A)$ 型。法向费米子没有零模，切向零模量子化给出 $\MC$ 上的旋量，即 $K_{\MC}^{1/2}\otimes(\bigoplus_j\Omega^{0,j})$；拓扑超荷 $Q$ 化为 $\bar\partial$，于是物理态恰是 $K_{\MC}^{1/2}$ 的全纯截面——这正好是 $\Dmod_{K_{\MC}^{1/2}}$ 的定义模（\src{B pp.197--198}）。第二个例子才是几何 Langlands 要用的：type F 膜，即卷在 Hitchin 纤维上、带平坦 Chan--Paton 线丛的膜。它是 $(A,B,A)$ 型，在 $K$ 中是 A-膜，所以按上面的命题自动给出 D-模；而它同时又是 zerobrane 的 S-对偶像，即 magnetic eigenbrane（\src{B p.198}）。

「magnetic eigenbrane」与「D-模」的对应，值得单独说清楚。特征膜的定义是：某条线算子作用在它上面时，只把它乘以一个固定的向量空间（或局部系统），而不是把它送到别的膜。在线算子一侧，这体现为「本征值」；翻译到代数一侧，同一个作用必须实现为该模层的一个自同态——因为 $(B_{c.c.},-)$ 把膜的作用映成模的作用。于是物理上的「本征值」就变成层论上的「Hecke 算子的本征作用」，而 magnetic eigenbrane 经翻译得到的那层，就是 Hecke 特征层。这一步与第 6 讲「't~Hooft 算子 $=$ Hecke 算子」是同一件事的两半：算子对应建立了，作用对象才第一次有意义（\src{B pp.184--185}）。

这里有一个必须承认的技术缝隙：层的粘接只对 $T^*V$ 型开集严格成立，精确理论中的完全层化涉及 Moyal 乘积的收敛问题，原文自己也没有做完（\src{B p.190}）。所以「每个 A-膜给出一个 D-模」在数学上仍是一条待修补的构造，而不是定理。

还要把话说白：原文这一步的口吻是「物理上应该」——c.c.\ 膜、$(B_{c.c.},B')$ 弦、层化，每一步都是物理上自然的构造，但没有一条定理保证它真的造出一个 D-模；数学上严格的对应要等后来纯代数几何的工作，下一讲再说。

\begin{figure}[htbp]
\centering
\begin{tikzpicture}[font=\small,
  bx/.style={draw=kwblue!70,fill=kwblue!4,rounded corners=1.5pt,align=center,
             inner sep=4pt,text width=3.15cm,font=\footnotesize},
  arr/.style={-{Stealth[length=2mm]},thick}]
\node[bx,draw=kwred!70!black,fill=kwred!5] (cc) at (0,2.3) {c.c.\ 膜 $B_{c.c.}$\\space-filling，$(A,B,A)$};
\node[bx] (alg) at (5.4,2.3) {$(B_{c.c.},B_{c.c.})$ 弦\\$=\Dmod_{K_{\MC}^{1/2}}$};
\node[bx] (a) at (0,0) {任意 A-膜 $B'$\\（复结构 $K$ 中）};
\node[bx] (mod) at (5.4,0) {$(B_{c.c.},B')$ 弦\\$=$ 上述代数上的模};
\node[bx,draw=kwgreen!65!black,fill=kwgreen!5] (dm) at (10.9,0) {$\MC$ 上 D-模\\$\Psi=h$ 时为普通 D-模};
\draw[arr] (cc) -- (alg) node[midway,above,font=\scriptsize]{自同态};
\draw[arr] (a) -- (mod) node[midway,above,font=\scriptsize]{张成模};
\draw[arr] (alg) -- (mod);
\draw[arr] (mod) -- (dm) node[midway,above,font=\scriptsize]{层化};
\draw[arr,dashed,draw=kwred!75!black] (dm.south) to[bend left=22]
  node[below,font=\scriptsize,text=kwred!75!black]{$\Psi\neq0$：把 $K_M^{1/2}$ 换成 $K_M^{1/2}\otimes L^{f(\Psi)}$} (mod.south);
\end{tikzpicture}
\caption{翻译器的工作方式。c.c.\ 膜的自同态代数给出微分算子层；任意 A-膜通过 $(B_{c.c.},B')$ 弦自动成为该代数的模，于是得到一个 D-模。$\Psi\neq0$ 时线丛指数要按 $f(\Psi)$ 调整（\src{B p.196}；\src{B p.207}）。}
\label{fig:translator}
\end{figure}

\subsection{量子化参数 $f(\Psi)$：普通 D-模与扭化 D-模}

为什么不能只在 $\Psi=0$ 处做翻译？因为几何 Langlands 对偶本身是交换 $\Psi=0$ 与 $\Psi=\infty$ 的 S 变换，$S$ 作用在参数上是 $(G,\Psi)\mapsto(\Ldual G,-\Psi/n_g)$，$T$ 作用是 $\Psi\mapsto\Psi+1$；全 S-对偶群把所有有理 $\Psi$ 串起来（\src{B p.201}）。若翻译器只在一个点上工作，就无法陈述对偶本身。于是 \S11.3 必须把 c.c.\ 膜推广到任意 $\Psi$。

一上 $\Psi\neq0$ 就遇到新障碍：A-模型的 B-场变成 $B=-\omega_I\Re\tau$，它沿 Hitchin 纤维不为零也不平坦，而带 B-场的 A-膜条件 $F+B|_Y=0$ 要求 $[B|_Y/2\pi]$ 是整上同调类；$\theta\neq2\pi n$ 时 Hitchin 纤维上根本没有合法的 rank-one A-膜（\src{B pp.202--203}）。原文的办法是改造探针：空间充满膜要解
\begin{equation}
\bigl(\omega^{-1}(F+B)\bigr)^2=-1,\qquad \omega=(\Im\tau)\omega_K ,
\label{eq:cc-psi}
\end{equation}
其解为 $F=\Im\tau\cos q\,\omega_J$、$\sin q=-\Re\tau/\Im\tau$（\src{B p.204}）。此时 $\MH$ 上的全纯结构不再是 $I$，而是 $I(\Psi)=I\cos q-J\sin q$；纤维方向的结构也不再是普通联络，而是 $\lambda$-连接，$\lambda=\ii\tan(q/2)$（\src{B p.204}）。换句话说，$\Psi$ 不只改变一个系数，它改变了「哪个复结构是拓扑的」。

代数结构随之改变：c.c.\ 弦层从 $\Dmod_{K_M^{1/2}}$ 变成作用在 $K_M^{1/2}\otimes L^{\Psi}$ 上的微分算子层（\src{B p.203}）。原文接着用整数点把它钉死：$\Psi\to\Psi+n$ 使 $B/2\pi$ 平移 $c_1(L^{-n})$，于是在 $\Psi=n$ 处可以取支撑在 $\varphi=0$、Chan--Paton 丛为 $L^n$ 的膜 $B_n$，量子化 $(B_{c.c.}^n,B_n)$ 弦得到的模恰是 $\Dmod_{K^{1/2}\otimes L^n}$ 的模（\src{B p.205}）。从这里出发，原文先假设 $f$ 线性，由整数点 $f(n)=n$ 得 $f(\Psi)=\Psi$；再用欧氏作用量证明线性：把它写成「$T^*\MC$ 上恰当项 $+$ 从 $\MC$ 拉回的闭二形式」，后者的系数正比于 $\Re\tau$，而粘接修正只需算到 $1/(\Im\tau)^2$ 阶，在这一阶上严格线性（\src{B pp.206--207}）。最后一个常数来自 $\MC$ 的典范线丛同构于 $L^{-2h}$（$h$ 为 $G$ 的对偶 Coxeter 数），于是任意 $\Psi$ 处的代数层是 $L^{\Psi-h}$ 上的全纯微分算子层；取 $\Psi=h$ 就得到\textbf{普通}微分算子层，即作用在函数上的 D-模（\src{B p.207}）。

所以本讲给出的具体式子是
\begin{equation}
f(\Psi)=\Psi ,
\label{eq:f-Psi}
\end{equation}
来源是原文 \S11.3 的两步论证（整数点定值 $+$ 微扰线性），把典范丛的常数 $\!-h$ 吸收进去之后，普通 D-模落在 $\Psi=h$。这里的不确定度要如实标注：$f(\Psi)=\Psi$ 的推导是微扰的，只算到 $1/(\Im\tau)^2$ 阶，线性性依赖「$\Psi$ 是唯一参数」这一物理假设，不是非微扰构造；整数 $\Psi$ 处的值是硬锚点，非整数 $\Psi$ 处的结论则软得多（\src{B p.205}；\src{B p.207}）。

最后纠一个常见的记混。$\Psi=\ii$ 确实是一个特殊值，但那是 \S5 镜像对称的特殊值：在 $\Psi=\ii$ 上，对 $J$ 的 B-模型等价于对 $K$ 的 A-模型（\src{B p.71}）。它管的是「哪个模型描述同一个物理」，与 D-模的「通常／扭化」不是一回事。在 \S11.3 的归一化下，普通 D-模出现在 $\Psi=h$；若换用某些 D-模文献把普通 D-模记为扭化参数为零的约定，则此处的扭化参数应读作 $\Psi-h$。把这个约定记住，就不容易在 $f(\Psi)$ 上翻车。

\begin{keybox}
c.c.\ 膜之所以能当翻译器，是因为它空间充满、且曲率由 $\omega_J$ 典范地定死，于是它的开弦代数恰是微分算子层 $\Dmod_{K_M^{1/2}}$，任何 $K$ 中的 A-膜自动成为它的模；把 $\Psi$ 从 $0$ 调开，只是把这个层扭成 $L^{\Psi-h}$，$f(\Psi)=\Psi$，在 $\Psi=h$ 处回到普通 D-模。
\end{keybox}

\begin{follow}
\begin{itemize}
\item 主线是 \S11.1--\S11.3（\src{B pp.185--207}）。先读 \S11.1 开头的余各向同性膜定义与 (11.1)，确认 $F$ 为什么由 $\omega_J$ 定死、$(a,b)$ 族里为什么只剩 $b=0$（\src{B pp.185--187}）。
\item 再读 \S11.2 开头两段与 $U(1)$ 例子（\src{B pp.196--201}），重点只有一句：为什么任意 A-膜都是 c.c.\ 代数的模。
\item \S11.3 可跳读：只抓 (11.30)--(11.34) 与最后三页（\src{B pp.204--207}），$\lambda$-连接的细节不影响对 $f(\Psi)$ 的理解。
\item 时间紧就先背 $\Psi$ 轴上的三个点：$\Psi=0$ 给 $\Dmod_{K^{1/2}}$，$\Psi=n$ 给 $\Dmod_{K^{1/2}\otimes L^n}$，$\Psi=h$ 给普通 D-模（\src{B pp.205--207}）。
\end{itemize}
\end{follow}

\begin{checkbox}
\begin{itemize}
\item 为什么普通 Lagrangian A-膜不能当翻译器？「空间充满」这一条在机制里具体用在哪一步？
\item 在 $\MH$ 上，$(11.1)$ 的一般解有一整族，为什么只有 $b=0$ 那一支能定义膜？「典范」的「典」在哪里？
\item 从开弦代数走到 $\Dmod_{K_{\MC}^{1/2}}$，哪一步用到时间反演？它把线丛指数 $L$ 钉成了什么？
\item $f(\Psi)=\Psi$ 的论证用了哪两个输入？哪一个是硬锚点，哪一个是微扰假设？不确定度在哪一页？
\item $\Psi=h$ 与 $\Psi=\ii$ 各自特殊在哪里？为什么不能把 $\Psi=\ii$ 说成「普通 D-模所在的位置」？
\end{itemize}
\end{checkbox}
